%% LyX 2.4.4 created this file.  For more info, see https://www.lyx.org/.
%% Do not edit unless you really know what you are doing.
\documentclass[english]{article}
\usepackage[T1]{fontenc}
\usepackage[latin9]{inputenc}
\usepackage{amsmath}
\usepackage{amsthm}
\usepackage{amssymb}
\usepackage{geometry}
\geometry{verbose,tmargin=1in,bmargin=1in,lmargin=1in,rmargin=1in}
\PassOptionsToPackage{normalem}{ulem}
\usepackage{ulem}

\makeatletter
%%%%%%%%%%%%%%%%%%%%%%%%%%%%%% Textclass specific LaTeX commands.
\numberwithin{equation}{section}
\numberwithin{figure}{section}
\newlength{\lyxlabelwidth}      % auxiliary length 

%%%%%%%%%%%%%%%%%%%%%%%%%%%%%% User specified LaTeX commands.
\usepackage{fancyhdr}
\usepackage{lastpage}
\pagestyle{fancy}
\usepackage{enumitem}
\usepackage{ifthen}
\everymath=\expandafter{\the\everymath\displaystyle}
%\DeclareMathSizes{10}{10}{10}{10}
\usepackage{multicol}

\usepackage{tikz}
\usetikzlibrary{patterns}
\usetikzlibrary{plotmarks}
\usepackage{pgfplots}
\pgfplotsset{compat=newest}

\usepackage{calc}
\usepackage[nomessages]{fp}

\usepackage{datetime}
% Added by lyx2lyx
\setlength{\parskip}{\smallskipamount}
\setlength{\parindent}{0pt}

\makeatother

\usepackage{babel}
\begin{document}
\renewcommand{\author}{Dr.\,James Ashe}

\newcommand{\classNum}{439}

\newcommand{\classSec}{A}

\newcommand{\nameType}{\hspace{2cm} Name: solutions}

\newcommand{\examType}{Homework 2}

\newcommand{\Date}{\formatdate{6}{2}{2014}} %%\AdvanceDate[12] \today

%\newdate{my_date}{\day}{\month}{\year}%   
%\AdvanceDate[-5] 
%\displaydate{my_date}

%%First page's header%%%%%%%%%%%%%%%%%%%%%%
\fancypagestyle{firststyle} %%header settings for first pagr
{
	\fancyhead[L]{MTH \classNum{}(\classSec{}) \author{}\\
	\textbf{\examType{}} \Date{}}
	\fancyhead[C]{\textbf{\nameType}}
	\setlength{\headsep}{14pt}
}
%%%%%%%%%%%%%%%%%%%%%%%%%%%%%%%%%%%%%%%%%%

%%Next page's header%%%%%%%%%%%%%%%%%%%%%%%
%\setlist{nolistsep} %eliminates space between list items
\lhead{MTH \classNum{}(\classSec{})} 
\chead{\textbf{\nameType}} 
\cfoot{\thepage{}~of~\pageref{LastPage}} 
\setlength{\headsep}{5pt} 
%%%%%%%%%%%%%%%%%%%%%%%%%%%%%%%%%%%%%%%%%%%

%%Various settings; see the preamble for other settings%%%%%
%%\setenumerate[0]{leftmargin=12pt,itemindent=0pt,labelwidth=0pt}
\renewcommand{\labelenumi}{\textbf{\arabic{enumi}.}}
\renewcommand{\labelenumii}{\textbf{\alph{enumii}.}}
\thispagestyle{firststyle}
%%%%%%%%%%%%%%%%%%%%%%%%%%%%%%%%%%%%%%%%%%

\newcommand{\xspace}{\vspace{0cm}}

\textbf{Homework 2 due: \formatdate{13}{2}{2014}}

\emph{Choose any 10 problems from the following exercises in chapter
2 of the notes.}

\textbf{2.1 }Suppose $f(x_{1})=f(x_{2})$ for any linear function
$f\colon\mathbb{R}\rightarrow\mathbb{R}$ with a nonzero, defined
slope. $f(x_{1})=f(x_{2})\Rightarrow mx_{1}+b=mx_{2}+b\Rightarrow x_{1}=x_{2}$.
So $f$ is one-to-one.\vfill\xspace

\textbf{2.2 }Let $b\in\mathbb{Z}$. $f(b+0.5)=\left\lfloor b+0.5\right\rfloor =b$.
So $f$ is onto. To see that $f$ is not one-to-one, $f(1.5)=1=f(1)$
but $1.5\neq1$.\vfill\xspace

\textbf{2.3 }Recall that modulo operator distributes over addition
and multiplication. So for any two integers $a$ and $b$ we have
the following:

$f(a+b)=(a+b)\bmod n=(a\bmod n+b\bmod n)\bmod n=a\bmod n+b\bmod n=f(a)+f(a)$.

$f(ab)=(ab)\bmod n=((a\bmod n)(b\bmod n))\bmod n=a\bmod nb\mod n=f(a)f(b)$. 

So $f$ is a homomorphism. \vfill\xspace

\textbf{2.4 $f(1.5+1.5)=f(3)=3\neq1+1=f(1.5)+f(1.5)$ }\vfill\xspace

\textbf{2.5 }This problem should have read ``determine if $f$ is
a homomorphism''.

$f$ is \uline{not} a homomorphism under the usual operations since
$f(2+3)=2^{2+3}=2^{2}2^{3}=32\neq12=2^{2}+2^{3}=f(2)+f(3)$. However
it is a homomorphism under addition and multiplication defined in
the range as follows: $2^{x}\oplus2^{y}=2^{x+y}$ and $2^{x}\otimes2^{y}=2^{xy}$.
$f(x)\oplus f(y)=2^{x+y}=f(x+y)$, and $f(x)\otimes f(y)=2^{xy}=f(xy)$.

\vfill\xspace

\textbf{2.6 }The key here is to look at how the identity is mapped. 

The identity of $\mathbb{Z}_{n}$ is $1$. Let $f\colon\mathbb{Z}_{n}\rightarrow\mathbb{Z}_{n}$
be a homomorphism, and suppose that $f(1)=a$ ($1$ must map to something).
Since homomorphisms are operation preserving $f(x)=f(\underset{x\mbox{ times}}{\underbrace{1+1+\cdots+1})}=f(1)+f(1)+\cdots+f(1)=xf(1)=xa$,
and $a=f(1)=f(1\cdot1)=f(1)f(1)=a\cdot a=a^{2}$ so $k=k^{2}$.

For the second part: $g(f(a+b))=g(f(a)+f(b))=g(f(a)+f(b))=g(f(a))+g(f(b))$,
and $g(f(ab))=g(f(a)f(b))=g(f(a))g(f(b))$. So yes $g\circ f$ is
a homomorphism.\vfill \xspace

\textbf{2.7 $2^{2}\neq2\bmod8$},\textbf{ }so it is not a ring homomorphism
by exercise 2.6. Alternatively, $f(3\cdot3)=9\cdot2=2\bmod8$, but
$f(3)f(3)=6\cdot6=4\bmod8$ so $f(3\cdot3)\neq f(3)f(3)$.\vfill \xspace

\textbf{2.8 }Yes, it is a ring homomorphism by exercise 2.6 since
$6^{2}=6\bmod30$. Alternatively, $f(a+b)=6(a+b)=6a+6b$, and $f(ab)=6(ab)=6\bmod30\cdot(ab)\bmod30=36\bmod30\cdot(ab)\bmod30=6a\cdot6b\bmod30=f(a)f(b)$.\vfill \xspace

\textbf{2.13 }Suppose that $f\colon\mathbb{Q}\rightarrow\mathbb{Z}$
is an isomorphism. Then $f$ is homomorphic, and $f(1)=f(\frac{1}{2}\cdot2)=f(\frac{1}{2})f(2)$.
Since $f$ is an isomorphism it must also be onto; so by theorem 2.1
(part 7) $f(1)=1$. This implies that $f(\frac{1}{2})f(2)=1$. $f$
maps to $\mathbb{Z}$ so $f(\frac{1}{2})=a$ and $f(2)=b$ for some
integers $a$ and $b$. But this implies that $ab=1$ which is only
true for integers $a=b=1$. A contradiction because then $f(1)=f(\frac{1}{2})=1$,
and $f$ is not one-to-one.

Similarly we can show that $\mathbb{R}$ is not isomorphic to $\mathbb{Z}$.

\vfill \xspace

\textbf{2.14 }$f$ cannot be an isomorphism because $f$ is not one-to-one.
$f(0)=f(n)=0\bmod n$. Alternatively, $|\mathbb{Z}|\neq|\mathbb{Z}_{n}|$
so $f$ cannot be one-to-one by the pigeon hole theorem. So there
is no isomorphism between them.

\vfill \xspace

\textbf{2.15 }Suppose that $m>n$ then $m=|\mathbb{Z}_{m}|>|\mathbb{Z}_{n}|=n$.
Then any function $f\colon\mathbb{Z}_{m}\rightarrow\mbox{\ensuremath{\mathbb{Z}}}_{n}$
cannot be one-to-one by the Pigeon hole principle. If $n<m$ then
$f$ cannot be onto, so $f$ can be an isomorphism only if $n=m$.

\vfill
\end{document}
